% Source: ACM_Templates.pdf pp. 40-66. User Node/tr/w/pushup skeleton, pp. 42-43.
\subsection{线段树：单点赋值与区间和}
\label{ds:seg-point}
\AlgorithmInfo
    {提供熟悉的节点合并骨架，维护单点赋值和区间和。}
    {建树 $O(n)$；修改、查询 $O(\log n)$；空间 $O(n)$。}
    {$n\ge1$；闭区间、1 下标；区间和及中间运算在 64 位范围内。}
\subsubsection*{原理}
每个节点对应一段连续区间，左右儿子各维护一半。\InlineCode{pushup} 合并两个儿子的区间和。修改只沿根到一个叶子的路径进行；查询将目标区间拆成 $O(\log n)$ 个节点，按从左到右的顺序合并。

保留 \InlineCode{pushup(Node&,const Node&,const Node&)} 与 \InlineCode{pushup(int)} 两层写法，方便以后替换信息；前者也用于合并查询得到的临时节点。

\subsubsection*{使用教程}
构造时指定 $n$，把初值放进 \InlineCode{w[1..n]}，再调用 \InlineCode{build(1,1,n)}。\InlineCode{modify(1,idx,val)} 将单点改为 $val$，\InlineCode{query(1,L,R)} 返回 \InlineCode{Node}，其中 \InlineCode{sum} 为区间和。
\begin{lstlisting}[style=mtycpp,numbers=none]
PointSegTree seg(4);
seg.w={0,2,1,3,5};
seg.build(1,1,4);
seg.modify(1,2,7);
long long ans=seg.query(1,2,4).sum;
\end{lstlisting}
\needspace{7\baselineskip}
\subsubsection*{模板代码}
\CodeFile{code/data-structures/seg_point.cpp}
